% 4ページ版（../paper_asj4/main.tex）を、ページ数の制約を外して書き足した版（2026-11-04）。追加: 気導の基準、疑似喉マイク、字母単位の誤り（script 93）、スペクトログラム（script 94）、計算量。削除: 有声マスクの列
% ビルド: cd results/paper_asj6 && tectonic -X compile main.tex
% 数値の出典: results/GENERATIVE_RESULTS_2026-09-30.md §6・§7（特に §7.12〜§7.19）、results/review_response/stats_taps.csv
% 注意: 学会の正式テンプレートではなく、余白（上下18 mm・左右22 mm・段間6 mm）を合わせた自作の書式。投稿時は学会配布のスタイルファイルに載せ替える
\documentclass[xelatex,ja=standard,a4paper,twocolumn,10pt]{bxjsarticle}
\geometry{top=18mm,bottom=18mm,left=22mm,right=22mm,columnsep=6mm}
\setjamainfont[BoldFont={Hiragino Sans W6}]{Hiragino Mincho ProN}
\setjasansfont[BoldFont={Hiragino Sans W6}]{Hiragino Sans W3}
\usepackage{amsmath,amssymb}
\usepackage{graphicx}
\usepackage{booktabs}
\usepackage[font=small,skip=3pt]{caption}
\usepackage{cite}

\pagestyle{empty}
\setlength{\textfloatsep}{10pt plus 2pt minus 2pt}
\setlength{\floatsep}{6pt plus 2pt minus 2pt}

\makeatletter
\renewcommand{\section}{\@startsection{section}{1}{\z@}{1.0ex plus .3ex}{.5ex}{\normalfont\large\sffamily\bfseries}}
\renewcommand{\subsection}{\@startsection{subsection}{2}{\z@}{.7ex plus .2ex}{.3ex}{\normalfont\normalsize\sffamily\bfseries}}
\makeatother

\newcommand{\figdir}{../figures}
\setCJKfamilyfont{hangul}{AppleMyungjo}
\newcommand{\ko}[1]{{\CJKfamily{hangul}#1}}

\begin{document}

\twocolumn[{%
  \begin{center}
    {\Large\sffamily\bfseries 喉マイク音声の認識には、どんな前処理が効くのか\\[2pt]
      ——事前学習済み生成モデルの適応とWhisper用音声強調の役割——\footnotemark}\\[8pt]
    {春日　裕次（明治大学　総合数理学部　先端メディアサイエンス学科）}
  \end{center}
  \vspace{6pt}
}]
\footnotetext{What preprocessing helps throat-microphone speech recognition? Roles of pretrained generative restoration and Whisper-oriented speech enhancement, by KASUGA, Yuji (Meiji University).}

%────────────────────────────────────────
\section{はじめに}
喉マイクは頸部の振動を捉えるため環境騒音に強いが、高域や一部の子音の情報が乏しい。
喉マイクと気導マイクの同時収録コーパスTAPS\cite{taps}の喉マイク音声をWhisper-small\cite{whisper}でそのまま認識すると、文字誤り率（CER）は0.447に達する。

前処理としての音声強調（SE）は有力な対策に見えるが、信号の品質を改善しても、処理による歪みのために音声認識（ASR）を悪化させうることが知られている\cite{ochiai2024}。
喉マイクでも、汎用の帯域拡張や声質変換（TAPSで追加学習したものを含む）がCERを悪化させたと報告されている（正解テキストではなく、気導音の認識結果を基準にCERを計算）\cite{yamanaka2026}。
一方、Hauretら\cite{hauret2025rt}は本稿と同じ生成モデルをVibraVoxの喉マイクで追加学習し、1種の音素認識器で誤りの低下を報告した。
また、凍結したASRの損失で前処理を学習する方法\cite{dissen2024}や、生成型と予測型の出力を融合する方法\cite{shi2024}も提案されている。
本稿では、喉マイクに対して前処理が認識を助けるのはどんな条件か、ASRの損失で学習したSEが生成モデルの出力に何を加えるのか、それがなぜ効くのかを、系統の異なる複数の認識器で調べる。
目的は融合方式の提案ではなく、改善が残る条件と、その改善に関わる成分の分析である。

本稿は、TAPSで9種、VibraVoxで7種の認識器を用いてこれを分析し、次の4点を示す。
(1) 事前学習済みの生成モデルを対象の喉マイクで適応すると、喉マイク音声で追加学習したWhisperを含む9種すべてでCERが下がる。同じ更新回数で乱数初期値から学習したモデルや、別の喉マイクで適応した生成モデルをそのまま使うと悪化した。
(2) 気導音に喉マイクの平均的な振幅特性と4~kHzの帯域制限をかけた疑似喉マイク音声は、本物の喉マイクほど認識を難しくしなかった。また、TAPSの一部の発話に限った探索的な分析では、生成モデルは字母（ハングルの子音字・母音字）の認識誤りを広く減らしたが、気導音との差に対して残った誤りの割合は、Whisper-smallで激音（強い息を伴う子音）が平音より大きかった。
(3) Whisperの損失で学習したSEの出力を生成モデルの出力と振幅で混ぜると（融合）、Whisper系で生成モデル単体よりさらに少しCERが下がる（以下、この追加の改善を上乗せと呼ぶ）。上乗せはWhisper-small・medium・turbo・FT Whisperでは3回の学習すべてで有意であった。baseは3回中1回のみ有意で、Qwen3はどの回も有意でなかった。最初に学習したモデル（以下、元の学習）では、上乗せは時間平均した固定のゲインでは再現しなかった。融合と生成モデル単体の対数振幅の差を、スペクトル包絡に当たる滑らかな成分（包絡）だけにして加えると、medium・turbo・FT Whisperで有意に改善し、その量は融合の改善量の約半分〜9割（点推定）に当たった。TAPSでは同じ包絡の成分がXLSRとZipformerを悪化させたが、VibraVoxではCE-SEと融合が仏語版のXLSRとMMSも改善し、系統の異なる認識器への影響の正負はデータセットで変わった。
(4) 今回の2つのデータセットの間では、一方で適応した生成モデルを他方にそのまま用いると、どちらの組み合わせでも悪化したが、同じ手順で適応し直せば、VibraVoxの7種の認識器で、生成モデル（4サンプル平均）で処理なしよりCERが70〜82\%下がり（7種とも有意）、融合で点推定ではさらに3〜15\%下がった（有意は6種で、baseは$p=0.055$）。

%────────────────────────────────────────
\section{方法}
\begin{figure}[t]
  \centering
  \includegraphics[width=0.92\columnwidth]{\figdir/method_overview_4p.png}
  \caption{手法の全体像。喉マイク音声を生成モデルとWhisper用SEの両方に通し、出力の振幅を融合して9種の認識器で評価する}
  \label{fig:overview}
\end{figure}

\subsection{生成モデルによる復元}
手法の全体像を図\ref{fig:overview}に示す。
NVIDIAが公開しているflow matching\cite{lipman2023}型の音声復元モデル\cite{ku2024}を用いる。
複素短時間フーリエ変換（STFT）を扱う24層のTransformer（4.3億パラメータ）で、英語約6万時間の音声で、マスクした区間を復元するタスクにより事前学習されている。
これを喉マイク音声を条件、同時収録の気導マイク音声を目標として追加学習する（以下、適応。ASRの損失は用いない）。喉マイクと気導マイクの同時収録の組を対データと呼ぶ。
推論では標準正規分布の雑音から出発し、flow matchingの常微分方程式をEuler法20ステップで解いて生成する。
生成は出発点の雑音に依存してばらつくので、$N=4$個のサンプルのSTFT振幅を平均し、位相は1つ目のサンプルのものを用いる（以下、生成×4）。

\subsection{Whisper用SEと振幅融合}
TAPS論文\cite{taps}が公開しているSE-Conformer\cite{seconformer}（1210万パラメータ、以下TAPS SE。TAPSのtrainで、喉マイク音声から気導マイク音声への変換として再構成損失で学習済み）を初期値とし、交差エントロピー（CE）を含む次の損失で追加学習する（以下CE-SE）。
\begin{equation}
  L = L_\mathrm{recon}(\hat{s}, s_\mathrm{air}) + \lambda\, \mathrm{CE}(\mathrm{Whisper}(\hat{s}), y)
  \label{eq:loss}
\end{equation}
$\hat{s}$はSE出力、$s_\mathrm{air}$は気導音声、$y$は正解テキスト、$\lambda$はCE損失の重みである。CEは、$\hat{s}$からWhisperと同じ対数メルを微分可能な形で求めてWhisper-smallに入力し、正解テキストのトークン列をデコーダに与えたとき（言語は韓国語に指定）の、各トークンの予測の交差エントロピーである。
$L_\mathrm{recon}$は再構成損失で、TAPS SEの学習\cite{taps}と同じく、波形のL1損失と多重解像度STFT損失\cite{pwg}の和とした。後者は3つの解像度（FFT長512・1024・2048点）で、振幅のspectral convergenceと対数振幅のL1誤差を0.5ずつ足して平均したものである。
Whisper-smallの重みは更新せず、CE損失の勾配だけをWhisperを通してSEに伝える。
生成×4とCE-SEの出力のSTFT（512点、シフト128点、Hann窓）の複素係数をそれぞれ$G$、$C$とし、振幅を次式で融合して、位相は$G$のものを用いる（以下、この方法を単に融合と呼ぶ。$Y$は融合後の複素係数、$j$は虚数単位）。
\begin{equation}
  Y = \bigl((1-w)\,|G| + w\,|C|\bigr)\, e^{j\angle G}
  \label{eq:fusion}
\end{equation}
$\lambda=10$は$\{0, 0.1, 0.5, 1, 2, 5, 10\}$から、式\eqref{eq:fusion}のCE-SE側の重み$w=0.5$は生成1回の出力との融合で$\{0.25, 0.5, 0.75\}$から、いずれもdevでWhisper-smallのCERのみを見て選んだ。

%────────────────────────────────────────
\section{実験条件}
TAPS\cite{taps}（韓国語）の話者独立な分割（train 40話者・dev 10話者・test 10話者、各100発話）を用いた。音声はすべて16~kHzで扱った。TAPSの喉マイクは8~kHzで収録されており、4~kHz以上の成分は収録の時点から含まれない\cite{taps}。
生成モデルはtrainの4,000対で2万回の更新（1バッチは6.14秒の切り出し8個、学習率$10^{-4}$）だけ追加学習し、最終ステップの重みを用いた。
CE-SEは15秒を超える発話を除き、バッチ4発話、学習率$3\times10^{-4}$で最大50エポック学習し、devでの式\eqref{eq:loss}の損失が5エポック続けて下がらなければ止め、その最小のエポックの重みを用いた。
testは話者が学習と重ならないが、収録機器・言語は同じである。
融合を採用するかは、Whisper-small以外の認識器のtestの結果も見て決めたため、TAPS上の分析は探索的である。

認識器は、Whisperと同形式の対数メルを入力とする6種（本稿では便宜上Whisper系と呼ぶ。Whisper-small（損失と選択に使用）・base・medium・large-v3-turbo、TAPS trainの喉マイク音声と正解テキストで全パラメータを追加学習（fine-tuning。学習率$10^{-5}$、バッチ16、devのCERが最小のエポック）したWhisper-small（以下FT Whisper）、Qwen3-ASR-1.7B\cite{qwen3asr}）と、系統の異なる3種（波形を入力とし、CTC（connectionist temporal classification）で学習されたMMS-1B\cite{mms}（韓国語のアダプタ）とXLSR-53\cite{xlsr}を韓国語でfine-tuningした公開モデル（XLSR）、韓国語の自発音声コーパスKsponSpeechで学習された公開のZipformer\cite{zipformer} transducer）である。
Whisper系は言語を指定し、Whisperはビーム幅5で、CTCの2種とZipformerは貪欲法で復号した。
CERは文字（韓国語はハングルの音節1文字）単位で、句読点を除き（仏語は小文字化も）、空白も1文字として数え、発話ごとのCERを1.0で打ち切ってから全発話で平均した。全発話の編集距離の合計を、正解の文字数の合計で割る通常のCERでも、主な傾向は同じであった。Zipformerは空白を出力しないため、この数え方ではCERが高く出るが、空白を除いたCERでも主な条件の大小関係は同じであった。
VibraVoxでは、韓国語のみのFT WhisperとZipformerを除き、XLSRとMMSを仏語版に替えた7種（Whisper系5種、XLSR、MMS）を用いた。
話者ごとの平均CERで両側のWilcoxon符号順位検定を行い、条件の対ごとに、その対を評価したすべての認識器のp値をまとめてHolm法で補正した（全探索に対する補正ではない）。補正後の$p<0.05$を有意とした。TAPSの10話者では補正前のpは最小でも0.002であり、9種で補正した後の最小値は0.018である。
比較手法のSSL-MSEは、Satoら\cite{sato2025}の自己教師あり学習（SSL）表現の損失を参考にした対照で、WavLM-Large\cite{wavlm}の全層で表現の平均二乗誤差（MSE）を求めて平均し、重み0.3で再構成損失に加えてTAPS SEを追加学習した（重み0.3は、SSLの項が再構成損失の約13〜14倍になるように決めた。CE-SEの$\lambda=10$でのCEの項と同じ比である。学習は最大15エポックで、devの損失で選んだ）。原典は後半の層を重み付きで平均した表現の誤差を用いており、本稿の実装とは異なる。
CE-SEと生成モデルは、元の学習に加えて、乱数のシードと計算環境（数値精度を含む）を変えて2回ずつ学習し直し（学習し直し1・2）、同じ回どうしを融合して再現性を確かめた。評価は3回とも同じtestの10話者で行った。

%────────────────────────────────────────
\section{実験結果}
\subsection{前処理が効く条件}
表\ref{tab:main}に各手法のCERを示す。
生成×4は9種すべてでTAPS SEよりCERを下げた（$N=1$の1回の生成でもTAPS SEより15〜24\%低く、10話者全員で下がった）。
喉マイク音声で追加学習したFT Whisperでも、処理なしの0.138から生成×4で0.100、融合で0.092（処理なしより33\%低く、10話者全員で下がった）へ下がり、認識器を喉マイク音声で追加学習した後も前処理が効いた。
一方、同じ構造を同じ2万回の更新だけ乱数の初期値から学習したモデルはWhisper-smallで0.551、別の喉マイク（4.5節のVibraVox）で適応した生成モデルをそのまま用いると0.628（1回の生成）と、いずれも処理なし（0.447。VibraVoxで適応した生成モデルの評価と同時に認識し直した値は0.448）より悪化した。
融合後も、気導マイク音声の認識（Whisper-small 0.126）との差は残った。
今回比較した条件では、対象の喉マイクで適応した事前学習済みのモデルが有効であった。ただし、乱数初期値との比較は同じ2万回の更新に限ったもので、事前学習済みのモデルを適応せずに用いる条件は評価していない。このため、事前学習と適応のそれぞれが必要かは示しておらず、両者の寄与も分けていない。

TAPSの約10時間で4.3億パラメータのモデルを乱数から学習しきれていない可能性は残る（1210万パラメータのTAPS SEは同じデータで乱数から学習されている）。したがってこの対照が示すのは、同じ2万回の更新では、事前学習済みの重みから始めたときにだけ、この規模のモデルが約10時間の対データで有効になったことである。少ない対データでの適応の結果は4.5節に示す。

\begin{table}[t]
  \centering
  \caption{TAPS testのCER（太字は喉マイク入力の条件内での最小）。なし: 処理なし、TAPS: TAPS SE、W-FT: FT Whisper、気導: 同じ発話の気導マイク音声をそのまま認識した参考値（3種のみ測定）、*: 損失・選択に使用}
  \label{tab:main}
  \footnotesize
  \setlength{\tabcolsep}{3.0pt}
  \begin{tabular}{lcccccc}
    \toprule
    認識器 & なし & TAPS & CE-SE & 生成×4 & 融合 & 気導 \\
    \midrule
    W-small* & 0.447 & 0.230 & 0.196 & 0.169 & \textbf{0.163} & 0.126 \\
    W-base & 0.655 & 0.273 & 0.267 & 0.205 & \textbf{0.201} & -- \\
    W-medium & 0.337 & 0.196 & 0.169 & 0.150 & \textbf{0.144} & -- \\
    W-turbo & 0.257 & 0.170 & 0.141 & 0.131 & \textbf{0.125} & 0.103 \\
    W-FT & 0.138 & 0.146 & 0.121 & 0.100 & \textbf{0.092} & -- \\
    Qwen3 & 0.243 & 0.136 & 0.114 & 0.097 & \textbf{0.094} & 0.072 \\
    \midrule
    XLSR & 0.584 & 0.222 & 0.279 & \textbf{0.161} & 0.163 & -- \\
    MMS-1B & 0.497 & 0.350 & 0.354 & \textbf{0.284} & 0.285 & -- \\
    Zipformer & 0.997 & 0.542 & 0.791 & \textbf{0.433} & 0.483 & -- \\
    \bottomrule
  \end{tabular}
\end{table}

\subsection{疑似喉マイクとの比較と字母ごとの認識誤り}
喉マイクの難しさが平均的な周波数特性と帯域の狭さで説明できるかを確かめるため、疑似喉マイク音声（以下、疑似喉マイク）を作った。TAPS trainの最大1,000対から、気導音のエネルギーが発話内の最大から40~dB以内のフレームについて、喉マイクと気導マイクのパワーを周波数ごとに合計し（$S_\mathrm{throat}(f)$、$S_\mathrm{air}(f)$）、TAPS testの気導音のSTFT（512点、シフト128点）に振幅比$H(f)=\sqrt{S_\mathrm{throat}(f)/S_\mathrm{air}(f)}$をかけ、4~kHz以上の係数を0にし、RMSをtrainの喉マイク音声のRMSの中央値にそろえた。
Whisper-smallのCERは疑似喉マイク・気導音・本物の喉マイクの順に0.130・0.126・0.448であった（turbo 0.103・0.103・0.258、Qwen3 0.077・0.072・0.243。本物の喉マイクの値は表\ref{tab:main}とは別の回の認識結果）。
Whisper-smallでは疑似喉マイクと気導音のCERの差は0.003で、今回の平均的な振幅特性と4~kHzの帯域制限による疑似喉マイクは、本物の喉マイクの認識の難しさを再現しなかった。

\begin{figure}[t]
  \centering
  \includegraphics[width=0.78\columnwidth]{\figdir/spectrogram_gen_fusion.png}
  \caption{TAPS testの1発話（話者p00）の冒頭4秒のスペクトログラム。上から喉マイク（処理なし）、生成×4、融合、気導マイク。破線は4~kHz}
  \label{fig:spec}
\end{figure}

図\ref{fig:spec}に1発話のスペクトログラムを示す。
生成×4は喉マイクに無い4~kHz以上の成分を補い、その時間的な配置は多くの区間で気導マイクと対応して見えた。喉マイクの無音区間にある小さな雑音は、生成×4では消えた。
融合と生成×4の違いはこの図では目立たず、両者の差は4.4節で対数振幅の差として調べた。

どの字母の認識誤りが減ったかを見るため、句読点と空白を除いた認識結果と正解を字母（初声・中声・終声）に分解して単位コストの編集距離で対応づけ、正解の字母のうち置換か脱落になったものの割合を、字母の種類ごとに全発話で合算して求めた（以下、字母誤り率。挿入は数えない。表\ref{tab:jamo}）。
これは正書法上の字母の認識誤りで、実際に発音された音の誤りとは一致しない。
正解は数をハングルで書くが認識結果は数字で出ることがあり、\ko{칠}（7）・\ko{팔}（8）・\ko{천}（千）などに含まれる激音の誤りが水増しされるため、各認識器について、6条件（処理なし・TAPS SE・CE-SE・生成×4・融合・気導）のどの認識結果にも数字とラテン文字が出なかった発話に限った（Whisper-small 309、turbo 340、Qwen3 450発話。いずれも10話者を含む）。
この選択は認識結果に依存するため、選んだ発話はtest全体を代表しない。
処理なしでは激音（\ko{ㅋㅌㅍㅊ}）の字母誤り率が0.423と最も高く、同じ調音位置の平音（\ko{ㄱㄷㅂㅈ}）の0.277を上回った。
Whisper-smallでは、生成×4はすべての種類で字母誤り率を処理なしの約4分の1以下に下げた。処理なしと気導の差のうち生成×4で残った割合$r=(e_\mathrm{gen}-e_\mathrm{air})/(e_\mathrm{none}-e_\mathrm{air})$（$e$は各条件の字母誤り率。話者ごとに求めて中央値をとる）は、激音で18\%、平音で10\%であった。
激音の方が残る割合が大きく、Whisper-smallとturboでは有意であった（ともに8/10話者。この比較について3種の認識器内でHolm補正したが、比較を結果から選んだことは補正していない）。Qwen3では有意でなかった。
ただし、対応づけを逆方向から求める、3種に共通の283発話に限る、全1,000発話で集計する、のいずれでも、Whisper-smallの差は有意のまま残ったが、turboの差は有意でなくなった。
摩擦音も残る割合が大きい傾向があったが、破裂音との差は有意でなかった。
音の分類は音声学の区分に従ったが、激音と平音の比較は結果を見てから選んだもので、この分析は探索的である。

\begin{table}[t]
  \centering
  \caption{字母誤り率（Whisper-small、6条件のどの認識結果にも数字・ラテン文字が出なかった309発話）。残り: 本文の$r$の話者ごとの中央値（10話者とも分母は正）。平音・激音は破裂音と破擦音の初声}
  \label{tab:jamo}
  \footnotesize
  \setlength{\tabcolsep}{3.0pt}
  \begin{tabular}{lrcccc}
    \toprule
    種類 & 字母数 & なし & 生成×4 & 気導 & 残り \\
    \midrule
    母音 & 14104 & 0.282 & 0.050 & 0.021 & 13\% \\
    鼻音（\ko{ㄴㅁ}） & 1151 & 0.321 & 0.045 & 0.017 & 9\% \\
    平音（\ko{ㄱㄷㅂㅈ}） & 5112 & 0.277 & 0.034 & 0.010 & 10\% \\
    激音（\ko{ㅋㅌㅍㅊ}） & 1067 & 0.423 & 0.094 & 0.023 & 18\% \\
    摩擦音（\ko{ㅅㅆㅎ}） & 2712 & 0.225 & 0.054 & 0.021 & 19\% \\
    \bottomrule
  \end{tabular}
\end{table}

\subsection{ASR損失は何を足すか}
融合は、Whisper系6種すべてで生成×4よりCERを下げた（$-1.8$〜$-8.0$\%）。
各学習回のCERと補正後のp値を表\ref{tab:stat}に示す。
small・medium・turbo・FTでは、元の学習を含む3回とも有意であった（補正後$p\le0.020$）。baseは2回目のみ有意で、Qwen3はどの回も有意でなかった。
CE-SE単体はTAPS SEに比べて、元の学習ではbase以外のWhisper系5種を有意に改善したが、XLSRとZipformerは悪化させた（TAPS SEに対して、元の学習で$+26$\%・$+46$\%、学習し直した2回では約$+11$〜$12$\%。処理なしに対してはどちらも改善）。

\begin{table*}[t]
  \centering
  \caption{主要な比較の各学習回のCERと、話者単位のWilcoxon検定のHolm補正後のp値（TAPS test、3回とも同じ10話者。太字は$p<0.05$）。融合のp: 生成×4との比較。包絡（4.4節、元の学習のみ）のp: 対生成×4と、融合との比較}
  \label{tab:stat}
  \footnotesize
  \setlength{\tabcolsep}{3.0pt}
  \begin{tabular}{l ccc ccc ccc ccc}
    \toprule
    & \multicolumn{3}{c}{元の学習} & \multicolumn{3}{c}{学習し直し1} & \multicolumn{3}{c}{学習し直し2} & \multicolumn{3}{c}{包絡（元の学習）} \\
    \cmidrule(lr){2-4}\cmidrule(lr){5-7}\cmidrule(lr){8-10}\cmidrule(lr){11-13}
    認識器 & 生成×4 & 融合 & p & 生成×4 & 融合 & p & 生成×4 & 融合 & p & CER & p（対生成×4） & p（対融合） \\
    \midrule
    W-small* & 0.169 & 0.163 & \textbf{0.020} & 0.169 & 0.162 & \textbf{0.018} & 0.169 & 0.162 & \textbf{0.018} & 0.166 & 0.078 & 0.059 \\
    W-base & 0.205 & 0.201 & 0.422 & 0.205 & 0.201 & \textbf{0.029} & 0.204 & 0.200 & 0.137 & 0.204 & 0.984 & 0.193 \\
    W-medium & 0.150 & 0.144 & \textbf{0.018} & 0.151 & 0.143 & \textbf{0.018} & 0.150 & 0.142 & \textbf{0.018} & 0.146 & \textbf{0.049} & 0.193 \\
    W-turbo & 0.131 & 0.125 & \textbf{0.018} & 0.132 & 0.125 & \textbf{0.018} & 0.132 & 0.124 & \textbf{0.018} & 0.126 & \textbf{0.018} & 0.232 \\
    W-FT & 0.100 & 0.092 & \textbf{0.018} & 0.100 & 0.091 & \textbf{0.018} & 0.099 & 0.091 & \textbf{0.018} & 0.096 & \textbf{0.027} & 0.068 \\
    Qwen3 & 0.097 & 0.094 & 0.422 & 0.100 & 0.097 & 0.480 & 0.099 & 0.099 & 1.000 & 0.097 & 1.000 & \textbf{0.031} \\
    \midrule
    XLSR & 0.161 & 0.163 & 0.422 & 0.163 & 0.164 & 0.480 & 0.162 & 0.163 & 1.000 & 0.168 & \textbf{0.027} & \textbf{0.018} \\
    MMS-1B & 0.284 & 0.285 & 0.625 & 0.285 & 0.283 & 0.336 & 0.284 & 0.284 & 1.000 & 0.287 & 0.252 & 0.109 \\
    Zipformer & 0.433 & 0.483 & \textbf{0.018} & 0.432 & 0.437 & 0.770 & 0.431 & 0.438 & 0.336 & 0.510 & \textbf{0.018} & \textbf{0.031} \\
    \bottomrule
  \end{tabular}
\end{table*}

上乗せが単に2つの出力を平均した効果かを確かめるため、CE損失を用いないSEを同じ方法で融合した（表\ref{tab:ctrl}）。
TAPS SEや再構成損失のみで追加学習したSE（$\lambda=0$）との融合は、Whisper系で生成×4から$-1.7$〜$+3.1$\%とほぼ改善せず、CE-SEとの融合のCERは、これら2つの融合より6種すべてで有意に低かった。
ただしSSL-MSEとの融合でも、turbo・FTでは生成×4より有意にCERが下がった。
CE損失による学習は上乗せを大きくするが、CE損失がなくても上乗せの一部は得られる。

\begin{table}[t]
  \centering
  \caption{生成×4に融合するSEを変えた対照（元の学習、CER）。+TAPS: TAPS SE、+$\lambda$=0: 再構成損失のみで学習したSE、+SSL: SSL-MSE、+CE: CE-SE（本稿の融合）との融合。baseは本文のみ}
  \label{tab:ctrl}
  \footnotesize
  \setlength{\tabcolsep}{3.0pt}
  \begin{tabular}{lccccc}
    \toprule
    認識器 & 生成×4 & +TAPS & +$\lambda$=0 & +SSL & +CE \\
    \midrule
    W-small* & 0.169 & 0.172 & 0.172 & 0.168 & \textbf{0.163} \\
    W-medium & 0.150 & 0.150 & 0.151 & 0.147 & \textbf{0.144} \\
    W-turbo & 0.131 & 0.130 & 0.129 & 0.128 & \textbf{0.125} \\
    W-FT & 0.100 & 0.099 & 0.098 & 0.095 & \textbf{0.092} \\
    Qwen3 & 0.097 & 0.099 & 0.100 & 0.096 & \textbf{0.094} \\
    \bottomrule
  \end{tabular}
\end{table}

\subsection{上乗せに関わる成分}
融合が生成×4から何を変えたかを、元の学習のモデルで切り分けた（表\ref{tab:why}）。
対数振幅の差$D=\log|Y|-\log|G|$（$Y$: 融合）を用いる。$G$と$Y$は2.2節と同じSTFT（Hann窓）で求め、ゼロ振幅を避けるため振幅に$10^{-7}$を足してから自然対数をとった。
フレームの種類は、気導音を使わず生成×4の出力から決めた。フレームのパワーが発話内の最大から40~dB以上低いものを無音とし、無音でなく2~kHz未満のパワーの割合が0.5を超えるものを有声らしいフレーム、残りを無声とした。
まず、$D$を無音以外のフレームで時間平均した周波数ごとの値（以下、平均ゲイン）を、生成×4の対数振幅に加えた。
devで求めたもの、有声らしい・無声・無音ごとのもの、同じ発話の融合から求めたもの（発話ごとのイコライザに当たる。以下、発話別EQ）のいずれも、Whisper系で生成×4から$\pm2$\%以内であった（すべて有意でない）。
検討した平均ゲインでは、上乗せを再現できなかった。

次に、ケプストラムの30次未満だけを残す操作を$P$、恒等写像を$I$として、$D$を包絡$PD$と残りの細部$(I-P)D$に分け、片方だけをフレームごとに生成×4の対数振幅に加えた（両方を加えると融合の対数振幅に戻る。位相は生成×4に固定）。
$P$は、各フレームの$D$（257点）を512点の実ケプストラムに逆フーリエ変換し、ケフレンシー0〜29点（16~kHzで約1.9~ms）とそれに対称な点だけを残してフーリエ変換し、実部をとる操作である。30は、ピッチ周期（2.5~ms以上）より短くなるよう結果を見る前に決めた。
細部だけでは、どの認識器にも有意な改善は認められなかった。
包絡だけではmedium・turbo・FTで生成×4より3〜4\%下がり（有意、表\ref{tab:stat}）、融合の改善量の約半分〜9割に当たった（点推定で、区間は求めていない）。この3種では包絡だけと融合の差は有意でなかったが、これは両者が同等であることを示すものではない。Qwen3では包絡だけでは改善せず、融合は包絡だけより有意に低かった。
有声らしいフレーム（全体の約83\%、気導音のpYIN\cite{pyin}の有声判定との一致率は無音以外で0.88）でだけ融合の出力を使い、他のフレームでは生成×4をそのまま使うと、small・medium・turbo・FT・Qwen3で生成×4より有意に下がった。同じマスクを発話ごとに乱数で選んだ量だけ時間方向に巡回シフトしてもmedium・turbo・FTで有意に下がり、両者の差はどの認識器でも有意でなかった（マスクは3フレームの移動平均でなめらかにした）。効果が特定の区間に集まる証拠は得られなかった。
融合は、4~kHz以上の対数メルの気導音との平均絶対差（高域のずれ）と、有声らしいフレームの4~kHz未満の対数メルの回帰直線の傾き（スペクトル傾斜）を気導音に近づける一方、線形予測（LPC）で推定したフォルマントの帯域幅を広げた。LPCは、気導音のpYINで有声と判定したフレームで、8~kHzに下げてプリエンファシス（係数0.97）をかけた信号の25~msのハミング窓（10~msシフト）に10次で当て、90〜3800~Hzで帯域幅600~Hz未満の極を周波数の低い順に3つとった。
ただし、Whisper-small・medium・FT・Qwen3では、発話ごとのこれらの量の変化とCERの変化のスピアマン順位相関係数$\rho$は小さかった（$|\rho|$は約0.08以下。Zipformerでは約0.2の量もあった）。

包絡の成分は、系統の異なる認識器には逆に働いた。
Zipformerは発話別EQだけで生成×4より3.6\%改善した（10話者全員）が、包絡だけを加えると$+17.7$\%と、融合（$+11.7$\%）以上に悪化した。
TAPSでは、複数のWhisperへの上乗せも、XLSR・Zipformerの悪化も、包絡だけを加えたときに現れた（XLSRは包絡だけで$+3.9$\%）。

\begin{table}[t]
  \centering
  \caption{上乗せの切り分け（元の学習、CER）。EQ: 発話別EQ、包絡／細部: 融合と生成×4の対数振幅の差の包絡／細部だけを加える}
  \label{tab:why}
  \footnotesize
  \setlength{\tabcolsep}{3.0pt}
  \begin{tabular}{lccccc}
    \toprule
    認識器 & 生成×4 & EQ & 包絡 & 細部 & 融合 \\
    \midrule
    W-small* & 0.169 & 0.170 & 0.166 & 0.169 & 0.163 \\
    W-medium & 0.150 & 0.151 & 0.146 & 0.151 & 0.144 \\
    W-turbo & 0.131 & 0.131 & 0.126 & 0.132 & 0.125 \\
    W-FT & 0.100 & 0.100 & 0.096 & 0.099 & 0.092 \\
    Qwen3 & 0.097 & 0.096 & 0.097 & 0.099 & 0.094 \\
    \midrule
    Zipformer & 0.433 & 0.417 & 0.510 & 0.444 & 0.483 \\
    \bottomrule
  \end{tabular}
\end{table}

\subsection{別の喉マイクでの評価}
仏語の体内伝導マイクコーパスVibraVox\cite{vibravox}のspeech\_clean（静かな環境での読み上げ）の喉マイク音声（test 3,064発話・21話者、話者はtrainと重ならない）で評価した（表\ref{tab:vbx}）。testの文のうちtrainにも現れるのは1文で（TAPSでは1,000文中6文）、文の記憶でCERが下がった可能性は小さい。同時収録のヘッドセット気導マイク音声のCERは、devでWhisper-small 0.032、turbo 0.016、Qwen3 0.014であった。
TAPSで学習した生成モデルとSEをそのまま用いると、いずれも処理なしより悪化した。音量と帯域をTAPSにそろえても、処理なしより悪いままであった。
逆向き（VibraVoxで適応した生成モデルをTAPSに用いる）でも悪化しており（4.1節）、今回の2つのデータセットの間では、一方で適応した生成モデルは、適応し直さずには他方で使えなかった。ただし言語と機器が同時に変わっており、両者の寄与は分けていない。
VibraVoxのtrain（20,981対）で同じ手順の適応をやり直すと（生成モデルは事前学習済みから、SEはTAPS SEを再構成損失のみでVibraVoxに合わせてから、仏語の正解テキストでWhisper-smallのCEを用いてCE-SEを学習。$\lambda$、$w$、$N$はTAPSの値に固定）、7種の認識器で、生成×4は処理なしより70〜82\%、融合は生成×4より点推定でさらに3〜15\%CERを下げ（21話者、Holm補正後、生成×4対処理なしは7種とも$p<0.001$、融合対生成×4はbaseのみ$p=0.055$で有意でなく、他の6種は$p\le0.007$）、処理なし、生成×4、融合の順にCERが下がる関係が再現した。
TAPSと異なり、CE-SEと融合は仏語のCTCで学習された認識器（XLSR・MMS）も改善し、他系統への影響はデータセットに依存した。
Whisper-small・1回の生成では、trainの1,000対（約1.3時間、1万回の更新）の適応でも、全データ（約26時間、2万回の更新）の改善量の約86\%が得られた（1試行）。

\begin{table}[t]
  \centering
  \caption{VibraVox testのCER（仏語、7種。太字は各列の最小）。上段: TAPSで学習したまま、下段: VibraVoxで適応（再構成のみ: 再構成損失だけでVibraVoxに合わせたTAPS SE）。W-s・W-b・W-m・W-t: Whisper-small・base・medium・turbo。XLSR・MMSは仏語版（--は未評価）}
  \label{tab:vbx}
  \footnotesize
  \setlength{\tabcolsep}{1.6pt}
  \begin{tabular}{lccccccc}
    \toprule
    条件 & W-s & W-b & W-m & W-t & Qwen3 & XLSR & MMS \\
    \midrule
    処理なし & 0.422 & 0.536 & 0.366 & 0.269 & 0.242 & 0.275 & 0.258 \\
    TAPS SE & 0.711 & 0.755 & 0.718 & 0.654 & 0.594 & 0.501 & 0.577 \\
    生成×4 & 0.559 & 0.594 & 0.542 & 0.497 & 0.454 & -- & -- \\
    \midrule
    再構成のみ & 0.256 & 0.313 & 0.212 & 0.163 & 0.154 & 0.163 & 0.187 \\
    CE-SE & 0.153 & 0.226 & 0.114 & 0.081 & 0.081 & 0.128 & 0.131 \\
    生成×4 & 0.087 & 0.123 & 0.067 & 0.053 & 0.046 & 0.063 & 0.076 \\
    融合 & \textbf{0.079} & \textbf{0.119} & \textbf{0.059} & \textbf{0.045} & \textbf{0.041} & \textbf{0.058} & \textbf{0.067} \\
    \bottomrule
  \end{tabular}
\end{table}

\subsection{計算量}
融合は4.3億パラメータのモデルを20ステップ・4サンプル実行するため重い。Euler法のステップ数とサンプル数を減らした結果を表\ref{tab:steps}に示す。
5ステップ・1サンプルでは、この測定条件で生成部分の推論時間が約1/10になった（1発話あたり約4.1秒から約0.41秒、平均発話長は約9秒）。
本稿の20×4に対する5×1の融合のCERの変化は、Whisper系で$-0.03$〜$+3.6$\%で、Holm補正後に有意に悪化したのはbase（$+2.3$\%、$p=0.041$）であった。系統の異なる認識器では、XLSR（$+4.3$\%、$p=0.018$）とZipformer（$+11.0$\%、$p=0.018$）が有意に悪化した。5×4でも、XLSR（$p=0.047$）とZipformer（$p=0.018$）は20×4より有意に悪化した。

\begin{table}[t]
  \centering
  \caption{融合のステップ数・サンプル数とCER（元の学習、test 1,000発話）。条件: Euler法のステップ数×サンプル数、時間: 同じ100発話での生成モデルの推論時間（RTX PRO 6000、モデルの読み込みを含む一括の時間）。p: 20×4と5×1の比較の、話者単位のWilcoxon検定のHolm補正後のp値（9種で補正）。W-s・W-b・W-m: Whisper-small・base・medium、Zip.: Zipformer}
  \label{tab:steps}
  \footnotesize
  \setlength{\tabcolsep}{1.6pt}
  \begin{tabular}{lccccccc}
    \toprule
    条件 & 時間 & W-s & W-b & W-m & W-FT & XLSR & Zip. \\
    \midrule
    20×4（本稿） & 405 s & 0.163 & 0.201 & 0.144 & 0.092 & 0.163 & 0.483 \\
    5×4 & 115 s & 0.162 & 0.202 & 0.142 & 0.091 & 0.167 & 0.522 \\
    5×1 & 41 s & 0.164 & 0.206 & 0.144 & 0.095 & 0.170 & 0.536 \\
    p & & 1.000 & \textbf{0.041} & 1.000 & 0.186 & \textbf{0.018} & \textbf{0.018} \\
    \bottomrule
  \end{tabular}
\end{table}

%────────────────────────────────────────
\section{考察と限界}
前処理が喉マイク音声の認識を助けるかは、条件で反転する。
今回の設定では、事前学習済みの生成モデルを対象の喉マイクで適応したときは認識器の系統によらず効き、同じ2万回の更新で乱数初期値から学習したときや、別の喉マイクで適応したままのときは害になった。先行研究\cite{yamanaka2026}でCERが悪化した手法との違いが事前学習の有無にあるかは、本稿では確かめていない。
喉マイクの認識の難しさは、今回の平均的な振幅特性と帯域制限による疑似喉マイクでは再現されず、時間とともに変わる、音によって異なる劣化が関わっている可能性がある。
Whisper-smallでの探索的な分析では、生成後も激音の字母の誤りが相対的に多く残った。候補の一つは、激音に伴う強い気息（声帯の振動を伴わない息の雑音）の手がかりが喉マイクで弱まることである。ただし本分析は認識文字列に基づき、発声開始時間（VOT）や後続母音の基本周波数、表記と発音の違いを調べていないため、原因は特定できない。
ASR損失で学習したSEは、その上にWhisper系向けの小さな上乗せを足すもので、元の学習では、検討した平均ゲインでは再現せず、融合と生成×4の差の包絡だけで、medium・turbo・FTの上乗せの約半分〜9割が得られた。
融合による変化は気導音への一様な接近では説明できず、近づく特徴と遠ざかる特徴があった。他系統の認識器への影響はデータセットによって正負が変わり、原因は特定できていない。

限界として、今回は喉マイクごとに対データで適応しており、見たことのない機器への汎化は示せていない。VibraVoxでの評価も、話者は学習と重ならないが同一コーパス内の評価である。
評価は両データセットの収録音声のみで、喉マイクを使う主な理由である騒音下での効果は調べていない。
生成部分の処理時間は、4.6節の測定条件では約1/10まで減らせたが（モデルの読み込みを含む一括の時間で、逐次処理の遅延は測っていない）、ステップ数を減らすとXLSRとZipformerのCERは悪化した。
また、元の学習と学習し直した学習とでは計算環境も異なり、他系統への悪化の大きさは学習ごとに変わった。4.4節の分解は元の学習についての結果である。
用いた生成モデルの重みは非商用のライセンスである。

%────────────────────────────────────────
\section{まとめ}
喉マイク音声の前処理が認識を助ける条件と、上乗せに関わる成分を、TAPSで9種、VibraVoxで7種の認識器を用いて調べた。
事前学習済みの生成モデルを対象の喉マイクで適応すると、FT Whisperを含む9種すべての認識器でCERが下がった。
今回の平均的な振幅特性と帯域制限による疑似喉マイクは、喉マイクの認識の難しさを再現しなかった。探索的な分析では、Whisper-smallで生成後も激音の字母の誤りが相対的に多く残った。
Whisperの損失で学習したSEとの融合は、Whisper-small・medium・turbo・FT Whisperに、3回の学習すべてで有意な小さな上乗せを足した（baseは3回中1回のみ有意、Qwen3はどの回も有意でない）。元の学習では、融合と生成×4の差の包絡だけで、medium・turbo・FTの上乗せの約半分〜9割が得られた。
今回の2つのデータセットの間では、一方で適応した生成モデルは適応し直さずには他方で使えず、少量の対データでの適応と、見たことのない機器への汎化が今後の課題である。

\footnotesize
\let\oldbibitem\bibitem
\renewcommand{\bibitem}{\setlength{\itemsep}{0pt}\setlength{\parsep}{0pt}\oldbibitem}
\bibliographystyle{IEEEtran}
\bibliography{refs}

\end{document}
